% BUDGET: 0.15 pages (included in document header, not counted separately)

\begin{abstract}
Modern AI training datasets are assembled from thousands of sources through processing, tokenization, and packing---transformations that destroy the mapping between individual training records and their original contributors. When a contributor requests removal, audits demand license traceability, or quality issues require tracing specific data, this provenance gap forces trainers into catastrophic file-level over-deletion. We present OriginBlame (\texttt{ob}), a record- and token-level data provenance system that propagates author identity through data processing pipelines via a three-layer content-addressable architecture. \texttt{ob} resolves contributor-level queries into precise record sets through deterministic SHA-256 lookups, without model access or post-hoc inference. On a Chinese Wikipedia benchmark (220,000 pages, 482,543 contributors), record-level provenance reduces over-deletion from 2.8$\times$--44.8$\times$ (file-level) to near-parity, while pipeline integration adds 1--21\% throughput overhead (1--2\% HuggingFace, 9--21\% Datatrove). Provenance-based forget sets (training records targeted for removal) integrate with existing unlearning algorithms on a 1.7B model, though a gap remains: provenance identifies source records, but trained models store transformed knowledge beyond individual records. Our demonstration provides an interactive platform for real-time provenance queries, three-level revocation, and over-deletion comparison.
\end{abstract}
